\chapter{総合演習}\label{章:総合演習}

ここでは，問題のすぐ前に使う公式を置かず，式を見て方針を選ぶ練習をする。計算を始める前に，求めるものが原始関数，定積分の値，関数の性質，極限，あるいは誤差付きの近似のどれであるかを確かめたい。同じ積分記号でも，それによって使いやすい手法は変わる。

基礎問題は一つか二つの基本操作を確実につなぐ問題，標準問題は区間や式の見方を選ぶ問題である。発展問題では複数の章の内容を使うが，新しい一般定理を調べることは求めない。微分した式を元へ戻す，平方の積分を考える，端点の極限を独立に調べるなど，本書で扱った操作から進められる。ヒントは問題の後にまとめ，解答は巻末に置いた。最初の一手を自分で考えてから，必要な番号だけを参照してほしい。

\section{基礎問題}\label{節:総合基礎}

\begin{enumerate}[label=\textbf{\thechapter.\arabic*},leftmargin=*]
 \item\label{問:総合一割り算}
 次の定積分を求めよ。
 \[
  \int_0^1\frac{x^3+x^2+2x+1}{1+x^2}\dd x.
 \]
 \item\label{問:総合二置換}
 次の定積分を求め，置換後の区間も示せ。
 \[
  \int_0^1x\sqrt{1+x^2}\dd x.
 \]
 \item\label{問:総合三部分積分}
 次の定積分を求めよ。
 \[
  \int_1^{\e}x\log x\dd x.
 \]
 \item\label{問:総合四動く端点}
 $x>0$に対して$F(x)=\int_x^{x^2}\sqrt{1+t^2}\dd t$とする。
 $F'(x)$を求め，曲線$y=F(x)$の$x=1$における接線を求めよ。
 \item\label{問:総合五相補重み}
 $a>0$とする。次の積分を$a$で表せ。
 \[
  \int_{-a}^{a}\frac{1+x^2}{1+\e^x}\dd x.
 \]
 \item\label{問:総合六面積体積}
 曲線$y=x^2$と直線$y=x$で囲まれる領域の面積を求めよ。
 また，この領域を$x$軸のまわりに一回転してできる立体の体積を求めよ。
\end{enumerate}

%
\section{標準問題}\label{節:総合標準}

\begin{enumerate}[label=\textbf{\thechapter.\arabic*},leftmargin=*,start=7]
 \item\label{問:総合七区間の最適化}
 実数$a$に対して
 \[
  H(a)=\int_a^{a+1}(x^2-2x)^2\dd x
 \]
 とする。$H(a)$の最小値と，それを与える$a$をすべて求めよ。
 \item\label{問:総合八置換と部分積分}
 次の定積分を求めよ。
 \[
  \int_0^1x^3\log(1+x^2)\dd x.
 \]
 \item\label{問:総合九区分求積}
 次の極限を求めよ。
 \[
  \lim_{n\to\infty}\sum_{k=1}^{n}\frac{k^2}{n(n^2+k^2)}.
 \]
 \item\label{問:総合十軸を横切る回転}
 $0\leqq x\leqq1$，$x-1\leqq y\leqq x$で定まる領域を$x$軸のまわりに一回転してできる立体の体積を求めよ。
 各$x$での回転断面が円板になるか円環になるかを説明し，半径が切り替わる点を求めること。
 \item\label{問:総合十一漸化式}
 非負整数$n$に対して$K_n=\int_0^1\frac{x^n}{1+x}\dd x$とする。
 $K_0$と，$n\geqq1$で成り立つ漸化式を求め，$K_3$を計算せよ。
 さらに$K_n\to0$を示し，それを使って
 \[
  \lim_{n\to\infty}\sum_{k=1}^{n}\frac{(-1)^{k-1}}{k}=\log2
 \]
 を示せ。
 \item\label{問:総合十二両端判定}
 実数$\alpha$に対して，広義積分
 \[
  \int_0^{\infty}\frac{x^\alpha}{1+x^2}\dd x
 \]
 の収束条件を求めよ。また$\alpha=0$の場合の値を求めよ。
\end{enumerate}

%
\section{発展問題}\label{節:総合発展}

\begin{enumerate}[label=\textbf{\thechapter.\arabic*},leftmargin=*,start=13]
 \item\label{問:総合十三積分方程式}
 $[0,1]$で連続な関数$f$が，すべての$x\in[0,1]$について
 \[
  f(x)=1+\int_0^x(x-t)f(t)\dd t
 \]
 を満たすとする。$f$を求め，求めた関数が確かにこの等式を満たすことを確かめよ。
 微分方程式の一般的な解法は用いず，積の微分と基本定理から求めてよい。
 \item\label{問:総合十四二つの制約}
 $[0,1]$で連続な実数値関数$f$が
 \[
  \int_0^1f(x)\dd x=1,\qquad
  \int_0^1xf(x)\dd x=0
 \]
 を満たすとする。$\int_0^1f(x)^2\dd x$の最小値と，それを与える関数を求めよ。
 \item\label{問:総合十五媒介曲線}
 $0\leqq t\leqq\pi$で
 \[
  x=t-\frac12\sin2t,\qquad y=\sin^2t
 \]
 と表される曲線を考える。この曲線と$x$軸で囲まれる領域の面積，曲線の長さ，領域を$x$軸のまわりに一回転してできる立体の体積を求めよ。
 面積と体積の積分で，同じ部分を重複して数えていないことも確かめること。
 \item\label{問:総合十六近似と条件収束}
 次の二つの問いに，原始関数の初等的な表示を求めずに答えよ。
 \begin{enumerate}[label=(\arabic*),leftmargin=*]
  \item 有限テイラー公式を使い，$A=\int_0^1\sin(x^2)\dd x$を幅$\frac1{10000}$以下の有理数の区間に挟め。
  \item 広義積分$\int_1^{\infty}\sin(x^2)\dd x$が条件収束することを示せ。
 \end{enumerate}
 \item\label{問:総合十七ウォリス}
 非負整数$n$に対して
 \[
  C_n=\int_0^1(1-x^2)^n\dd x
 \]
 とする。$C_0$と$C_n$の漸化式を求め，$C_n$を有限の積で表せ。
 次にウォリスの積分との関係を使って，
 \[
  \lim_{n\to\infty}\sqrt n\,C_n,
  \qquad
  \lim_{n\to\infty}n^{\frac32}(C_n-C_{n+1})
 \]
 を求めよ。
 \item\label{問:総合十八相補と半角}
 $n$を正の整数とする。次の広義積分を，$n$を含む有限の積で表せ。
 \[
  J_n=\int_{-\pi}^{\pi}
  \frac{\sin^{2n}x}{(1+\e^x)\sqrt{1+\cos x}}\dd x.
 \]
 最初に両端での収束を確かめること。対称性，半角公式，置換積分，部分積分を組み合わせれば，ベータ関数やガンマ関数の公式を使わずに求められる。
\end{enumerate}

\clearpage
%
\section{ヒント}\label{節:総合ヒント}

以下では，計算の結論ではなく，最初の変形や見るべき条件を示す。解き直すときには，その変形を選んだ理由も言葉にしてみよう。

\begin{enumerate}[label=\textbf{\thechapter.\arabic*},leftmargin=*]
 \item 分子を$1+x^2$で割ると，余りの項が分母の微分と結び付く。
 \item 根号の中を$u=1+x^2$と置く。$x\dd x$をそのまま$\dd u$で表す。
 \item $\log x$を微分する側に選ぶと，残る積分は整式になる。
 \item $G(u)=\int_0^u\sqrt{1+t^2}\dd t$と置いて，$F(x)=G(x^2)-G(x)$と書く。
 \item $x=-t$で書き直した積分と元の積分を足す。$1+x^2$は偶関数である。
 \item 交点を求めた後で，上下関係と，回転軸からの二つの距離を調べる。
 \item 端点を微分すると$H'(a)$が得られる。値の計算には$y=a-\frac12$と置き，積分区間の中心から測ると整理しやすい。
 \item $u=1+x^2$と置けば，$x^3\dd x=\frac12(u-1)\dd u$になる。その後で部分積分する。
 \item 一項を$\frac1n$と$\frac{k}{n}$だけを使って表す。積分に直したら，分子を分母と定数の差にする。
 \item 区間$[x-1,x]$は常に$0$を含む。回転後の半径は，両端の軸からの距離の大きい方である。
 \item $K_n+K_{n-1}$では分母が消える。極限には$0\leqq K_n\leqq\int_0^1x^n\dd x$を使う。
 \item $0$付近では$x^\alpha$，無限大では$x^{\alpha-2}$と比較し，二つの条件を同時に満たす範囲を求める。
 \item 積分を$x\int_0^xf(t)\dd t-\int_0^xtf(t)\dd t$に分けて微分する。$f''=f$が得られたら，$\e^{-x}(f+f')$の微分を考える。
 \item まず二つの制約を満たす一次関数$\ell(x)$を求める。$g=f-\ell$と置いたとき，$\int_0^1\ell g\dd x$がどうなるかを調べる。
 \item $x'(t)=2\sin^2t$を使う。弧長では$\sqrt{x'(t)^2+y'(t)^2}$を整理してから積分する。
 \item (1) $\sin u$の六次テイラー多項式では，偶数次の係数が$0$になる。七階微分の符号と大きさを調べる。(2) $u=x^2$と置く。収束には部分積分，絶対値の発散には長さ$\pi$の区間ごとの下からの評価を使う。
 \item $x=\sin t$によって$C_n$を$\cos$の奇数べきの積分にする。隣り合うウォリス積分の比と積を組み合わせる。
 \item 指数関数を含む重みは，$x$を$-x$にすると相補的になる。その後，$x=2t$，さらに$u=\sin t$と置く。残る積分の$(1-u^2)$の指数は整数なので，部分積分で一つずつ下げられる。
\end{enumerate}

\smallskip\noindent 解答は\pageref{解答:第12章}ページ。
%
